\documentclass[aspectratio=169,professionalfonts]{beamer}

\usepackage{fontspec}
\usepackage{polyglossia}
\setmainlanguage{russian}
\setotherlanguage{english}
\defaultfontfeatures{Ligatures=TeX}
\setsansfont{Segoe UI}
\setmainfont{Segoe UI}
\setmonofont{Consolas}

\usepackage{tikz}
\usetikzlibrary{arrows.meta,positioning,calc,fit,shapes.geometric,backgrounds}
\usepackage{graphicx}
\usepackage{booktabs}
\usepackage{array}
\usepackage{multicol}

\definecolor{snowbg}{HTML}{0B0B0B}
\definecolor{snowpanel}{HTML}{141414}
\definecolor{snowpanelalt}{HTML}{1C1C1C}
\definecolor{snowline}{HTML}{EAEAEA}
\definecolor{snowmuted}{HTML}{BDBDBD}
\definecolor{snowaccent}{HTML}{FFFFFF}

\setbeamercolor{background canvas}{bg=snowbg}
\setbeamercolor{normal text}{fg=white,bg=snowbg}
\setbeamercolor{frametitle}{fg=white}
\setbeamercolor{title}{fg=white}
\setbeamercolor{subtitle}{fg=snowmuted}
\setbeamercolor{block title}{bg=snowpanelalt,fg=white}
\setbeamercolor{block body}{bg=snowpanel,fg=white}
\setbeamercolor{item}{fg=white}
\setbeamerfont{title}{series=\bfseries,size=\LARGE}
\setbeamerfont{frametitle}{series=\bfseries,size=\Large}
\setbeamerfont{normal text}{size=\small}

\setbeamertemplate{navigation symbols}{}
\setbeamertemplate{footline}{
  \leavevmode%
  \hbox{%
  \begin{beamercolorbox}[wd=.88\paperwidth,ht=2.5ex,dp=1ex,leftskip=1em]{author in head/foot}%
    \color{snowmuted}\insertshorttitle
  \end{beamercolorbox}%
  \begin{beamercolorbox}[wd=.12\paperwidth,ht=2.5ex,dp=1ex,rightskip=1em plus1fil]{date in head/foot}%
    \color{snowmuted}\insertframenumber/\inserttotalframenumber
  \end{beamercolorbox}}%
}

\newcommand{\spotlightbackground}{
\begin{tikzpicture}[remember picture,overlay]
  \fill[snowbg] (current page.south west) rectangle (current page.north east);
  \fill[white,opacity=0.025] (current page.south west) rectangle (current page.north east);
  \foreach \x in {0.03,0.08,...,0.97} {
    \draw[white,opacity=0.02,line width=0.2pt]
      ($(current page.south west)!\x!(current page.south east)$) --
      ($(current page.north west)!\x!(current page.north east)$);
  }
  \foreach \y in {0.04,0.1,...,0.96} {
    \draw[white,opacity=0.01,line width=0.2pt]
      ($(current page.south west)!\y!(current page.north west)$) --
      ($(current page.south east)!\y!(current page.north east)$);
  }
  \shade[inner color=white,outer color=white,opacity=0.12]
    ($(current page.north)+(0,-0.55)$) circle [radius=2.2cm];
  \shade[inner color=white,outer color=white,opacity=0.035]
    ($(current page.north)+(0,-1.2)$) circle [radius=5.7cm];
\end{tikzpicture}
}

\newcommand{\panel}[5]{
\begin{tikzpicture}
  \node[rounded corners=6pt,draw=white!18,fill=snowpanel,minimum width=#1,minimum height=#2,inner sep=0pt] (p) {};
  \node[anchor=north west,font=\bfseries\small,text=white] at ($(p.north west)+(0.22,-0.18)$) {#3};
  \node[anchor=north west,font=\fontsize{8}{10}\selectfont,text=snowmuted,align=left,text width=#4] at ($(p.north west)+(0.22,-0.55)$) {#5};
\end{tikzpicture}
}

\title[Snowdrop]{Snowdrop}
\subtitle{PostgreSQL-compatible acceleration and migration engine}
\author{}
\date{}

\begin{document}

\begin{frame}[plain]
\spotlightbackground
\vspace{1.2cm}
\begin{columns}[T,totalwidth=\textwidth]
  \begin{column}{0.62\textwidth}
    {\usebeamerfont{title}\color{white}Snowdrop}\par
    \vspace{0.3cm}
    {\large\color{snowmuted}Коммерческий взгляд на проект: \\ ускорение SQL-нагрузки без переписывания экосистемы PostgreSQL}\par
    \vspace{0.8cm}
    {\small\color{white}PostgreSQL wire protocol \quad $\cdot$ \quad RocksDB + Neo4j + columnar cache \quad $\cdot$ \quad migration workflow}
  \end{column}
  \begin{column}{0.38\textwidth}
    \begin{tikzpicture}[x=1cm,y=1cm]
      \node[rounded corners=8pt,draw=white!30,fill=snowpanel,minimum width=4.5cm,minimum height=4.8cm] {};
      \node[font=\bfseries\small,text=white] at (0,1.55) {Главная идея};
      \node[align=center,text width=3.8cm,text=snowmuted,font=\small] at (0,0.45) {
        Оставить клиентам привычный PostgreSQL-мир\\
        и ускорить те read-heavy сценарии,\\
        где монолитный row-store уже упирается в потолок
      };
      \node[draw=white,rounded corners=4pt,fill=snowpanelalt,minimum width=3.2cm,minimum height=0.75cm,text=white,font=\bfseries\small] at (0,-1.4) {0.8 ms on TPC-H Q6};
    \end{tikzpicture}
  \end{column}
\end{columns}
\end{frame}

\begin{frame}{1. Проблема рынка}
\spotlightbackground
\begin{columns}[T,totalwidth=\textwidth]
  \begin{column}{0.52\textwidth}
    \begin{tikzpicture}[node distance=0.55cm]
      \node[rounded corners=6pt,draw=white!20,fill=snowpanel,minimum width=6.4cm,minimum height=1.45cm,align=left,text width=5.7cm] (a) {
        \textbf{Вариант A: остаться на PostgreSQL}\\
        \color{snowmuted}Простая экосистема, но scan-heavy read path и reporting reads начинают дорого стоить
      };
      \node[rounded corners=6pt,draw=white!20,fill=snowpanel,minimum width=6.4cm,minimum height=1.45cm,below=of a,align=left,text width=5.7cm] (b) {
        \textbf{Вариант B: строить отдельный стек}\\
        \color{snowmuted}Усложнение архитектуры, ETL, дублирование данных, рост operational cost
      };
      \node[rounded corners=6pt,draw=white!35,fill=snowpanelalt,minimum width=6.4cm,minimum height=1.55cm,below=of b,align=left,text width=5.7cm] (c) {
        \textbf{Вариант C: Snowdrop}\\
        \color{snowmuted}Сохранить PostgreSQL-путь для клиентов и ускорить поддержанный SQL workload
      };
    \end{tikzpicture}
  \end{column}
  \begin{column}{0.48\textwidth}
    \begin{block}{Что именно болит}
      \begin{itemize}
        \item<1-> Operational dashboards прямо по production data
        \item<2-> Read-heavy backend endpoints с filters / aggregates / sorting
        \item<3-> Желание ускориться без полного rewrite архитектуры
        \item<4-> Нежелание терять SQL и драйверы PostgreSQL
      \end{itemize}
    \end{block}
  \end{column}
\end{columns}
\end{frame}

\begin{frame}{2. Что такое Snowdrop}
\spotlightbackground
\centering
\begin{tikzpicture}[node distance=0.55cm and 0.55cm, >=Latex,scale=0.96, transform shape]
  \node[rounded corners=8pt,draw=white!25,fill=snowpanel,minimum width=3.1cm,minimum height=1.1cm] (client) {\textbf{Client / ORM / psql}};
  \node[rounded corners=8pt,draw=white!35,fill=snowpanelalt,minimum width=3.3cm,minimum height=1.1cm,right=of client] (pg) {\textbf{PgWire endpoint}};
  \node[rounded corners=8pt,draw=white!35,fill=snowpanelalt,minimum width=3.6cm,minimum height=1.4cm,right=of pg] (orch) {\textbf{Orchestrator}\\\small parser + router + executor};

  \node[rounded corners=8pt,draw=white!20,fill=snowpanel,minimum width=3.1cm,minimum height=1.1cm,below=0.8cm of orch] (schema) {\textbf{Neo4j}\\\small schema metadata};
  \node[rounded corners=8pt,draw=white!20,fill=snowpanel,minimum width=3.1cm,minimum height=1.1cm,left=0.7cm of schema] (rocks) {\textbf{RocksDB}\\\small rows + secondary indexes};
  \node[rounded corners=8pt,draw=white!20,fill=snowpanel,minimum width=3.1cm,minimum height=1.1cm,right=0.7cm of schema] (col) {\textbf{Columnar cache}\\\small SIMD scan path};
  \node[rounded corners=8pt,draw=white!20,fill=snowpanel,minimum width=2.6cm,minimum height=1.1cm,below=0.7cm of schema] (wal) {\textbf{WAL}\\\small recovery};

  \draw[white,thick,->] (client) -- (pg);
  \draw[white,thick,->] (pg) -- (orch);
  \draw[white!70,thick,->] (orch) -- (rocks);
  \draw[white!70,thick,->] (orch) -- (schema);
\draw[white!70,thick,->] (orch) -- (col);
\draw[white!70,thick,->] (orch) -- (wal);
\end{tikzpicture}
\vspace{0.15cm}
{\color{snowmuted}\small\parbox{0.88\textwidth}{\centering Для приложения Snowdrop система выглядит как PostgreSQL-сервер, но внутри разделяет роли между schema metadata, durable row storage и accelerated read path.}\par}
\end{frame}

\begin{frame}{3. Почему архитектура устроена именно так}
\spotlightbackground
\begin{columns}[T,totalwidth=\textwidth]
  \begin{column}{0.33\textwidth}
    \panel{4.0cm}{2.3cm}{Neo4j}{3.3cm}{Хранит структуру базы: базы, таблицы, колонки, constraints, introspection metadata}
  \end{column}
  \begin{column}{0.33\textwidth}
    \panel{4.0cm}{2.3cm}{RocksDB}{3.3cm}{Хранит durable row data и secondary indexes без сетевого hop на каждый lookup}
  \end{column}
  \begin{column}{0.33\textwidth}
    \panel{4.0cm}{2.3cm}{Columnar cache}{3.3cm}{Даёт быстрый scan path для analytical predicates, aggregates и repeated reads}
  \end{column}
\end{columns}
\vspace{0.6cm}
\begin{block}{Ключевой принцип}
Один движок не обязан одинаково хорошо решать все задачи. Snowdrop разделяет обязанности так, чтобы критический read path не платил лишнюю цену за универсальность монолитного row-store.
\end{block}
\end{frame}

\begin{frame}{4. Как работает запрос}
\spotlightbackground
\begin{center}
\begin{tikzpicture}[node distance=0.9cm, >=Latex]
  \node[rounded corners=6pt,draw=white!25,fill=snowpanel,minimum width=2.8cm,minimum height=0.9cm] (q1) {SQL query};
  \node[rounded corners=6pt,draw=white!25,fill=snowpanel,minimum width=2.8cm,minimum height=0.9cm,right=of q1] (q2) {Parse + route};
  \node[rounded corners=6pt,draw=white!25,fill=snowpanel,minimum width=2.8cm,minimum height=0.9cm,right=of q2] (q3) {Choose path};
  \node[rounded corners=6pt,draw=white!25,fill=snowpanel,minimum width=2.8cm,minimum height=0.9cm,right=of q3] (q4) {Result};
  \draw[white,->,thick] (q1) -- (q2);
  \draw[white,->,thick] (q2) -- (q3);
  \draw[white,->,thick] (q3) -- (q4);

  \node[rounded corners=6pt,draw=white!18,fill=snowpanelalt,minimum width=2.8cm,minimum height=0.85cm,below left=1.0cm and -0.15cm of q3] (p1) {\small Indexed lookup};
  \node[rounded corners=6pt,draw=white!18,fill=snowpanelalt,minimum width=2.8cm,minimum height=0.85cm,below=1.0cm of q3] (p2) {\small Columnar SIMD scan};
  \node[rounded corners=6pt,draw=white!18,fill=snowpanelalt,minimum width=2.8cm,minimum height=0.85cm,below right=1.0cm and -0.15cm of q3] (p3) {\small Relation pipeline};
  \draw[white!70,->] (q3) -- (p1);
  \draw[white!70,->] (q3) -- (p2);
  \draw[white!70,->] (q3) -- (p3);
  \draw[white!55,->] (p1.east) to[out=0,in=-140] (q4.south);
  \draw[white!55,->] (p2) -- (q4.south);
  \draw[white!55,->] (p3.west) to[out=180,in=-40] (q4.south);
\end{tikzpicture}
\end{center}
\vspace{0.35cm}
\begin{itemize}
  \item<1-> Indexed path: selective lookup through RocksDB secondary indexes
  \item<2-> Columnar path: scan-heavy analytics with SIMD filtering
  \item<3-> Relation pipeline: joins, subqueries, set ops, windows, migration-safe subset
\end{itemize}
\end{frame}

\begin{frame}{5. Как работает миграция}
\spotlightbackground
\begin{center}
\begin{tikzpicture}[node distance=0.55cm, >=Latex]
  \node[rounded corners=6pt,draw=white!25,fill=snowpanel,minimum width=2.7cm,minimum height=0.9cm] (a) {1. Discover source};
  \node[rounded corners=6pt,draw=white!25,fill=snowpanel,minimum width=2.4cm,minimum height=0.9cm,right=0.45cm of a] (b) {2. Preflight};
  \node[rounded corners=6pt,draw=white!25,fill=snowpanel,minimum width=2.8cm,minimum height=0.9cm,right=0.45cm of b] (c) {3. Safe translation};
  \node[rounded corners=6pt,draw=white!25,fill=snowpanel,minimum width=2.4cm,minimum height=0.9cm,below=0.65cm of b] (d) {4. Load data};
  \node[rounded corners=6pt,draw=white!25,fill=snowpanel,minimum width=2.1cm,minimum height=0.9cm,right=0.45cm of d] (e) {5. Verify};
  \node[rounded corners=6pt,draw=white!25,fill=snowpanel,minimum width=2.1cm,minimum height=0.9cm,right=0.45cm of e] (f) {6. Report};

  \draw[white,->,thick] (a) -- (b);
  \draw[white,->,thick] (b) -- (c);
  \draw[white,->,thick] (c.south) -- ++(0,-0.35) -| (d.north);
  \draw[white,->,thick] (d) -- (e);
  \draw[white,->,thick] (e) -- (f);
  \node[text=snowmuted,align=left,font=\small] at ($(a.south)+(0,-1.7)$) {
    \texttt{migration-preflight}\\
    \texttt{migration-run}\\
    \texttt{migration-auto}
  };
  \node[draw=white!18,rounded corners=6pt,fill=snowpanelalt,minimum width=4.7cm,minimum height=1.3cm,align=left,text width=4.1cm] at ($(f.south)+(0,-1.45)$) {
    \textbf{Важно:}\\
    unsafe blockers не замалчиваются,\\
    а режутся рано и явно
  };
\end{tikzpicture}
\end{center}
\end{frame}

\begin{frame}{6. Метрики: где Snowdrop уже силён}
\spotlightbackground
\begin{columns}[T,totalwidth=\textwidth]
  \begin{column}{0.58\textwidth}
    \begin{tabular}{>{\raggedright\arraybackslash}p{2.4cm} >{\raggedleft\arraybackslash}p{1.5cm} >{\raggedleft\arraybackslash}p{1.6cm} >{\raggedleft\arraybackslash}p{1.8cm}}
      \toprule
      Query & Snowdrop & PostgreSQL & Signal \\
      \midrule
      Q6 & \textbf{0.8 ms} & 1.5 ms & 1.9x faster \\
      Q1 & 10.9 ms & \textbf{4.6 ms} & PG stronger \\
      Q14 & 26.1 ms & \textbf{1.5 ms} & JOIN gap \\
      Q3 & 44.3 ms & \textbf{1.6 ms} & JOIN gap \\
      \bottomrule
    \end{tabular}
  \end{column}
  \begin{column}{0.42\textwidth}
    \begin{block}{Как это правильно продавать}
      \begin{itemize}
        \item<1-> Не обещать ускорение всего подряд
        \item<2-> Показывать сильный сигнал на scan-heavy read path
        \item<3-> Честно говорить, что JOIN-heavy SQL пока слабее PostgreSQL
      \end{itemize}
    \end{block}
  \end{column}
\end{columns}
\vspace{0.35cm}
\color{snowmuted}\small TPC-H scale factor 1, warm-cache latencies. Snowdrop already wins on the workload class it was designed to accelerate.
\end{frame}

\begin{frame}{7. Где продукт реально полезен}
\spotlightbackground
\begin{columns}[T,totalwidth=\textwidth]
  \begin{column}{0.55\textwidth}
    \begin{block}{Хороший migration candidate}
      \begin{itemize}
        \item Ticketing / events platform
        \item Order management / e-commerce backend
        \item Billing / subscriptions
        \item CRM-like operational database
        \item Read-heavy service database
      \end{itemize}
    \end{block}
  \end{column}
  \begin{column}{0.45\textwidth}
    \begin{tikzpicture}
      \node[rounded corners=8pt,draw=white!20,fill=snowpanel,minimum width=4.8cm,minimum height=4.3cm] {};
      \node[font=\bfseries\small,text=white] at (0,1.55) {Типичный сценарий};
      \node[align=left,text width=4.0cm,text=snowmuted,font=\small] at (0,-0.1) {
        PostgreSQL-база большого backend-продукта\\
        + heavy reads\\
        + dashboards\\
        + aggregates\\
        + filters / sorts / pagination\\
        + желание ускориться без rewrite
      };
    \end{tikzpicture}
  \end{column}
\end{columns}
\end{frame}

\begin{frame}{8. Пример продукта}
\spotlightbackground
\begin{columns}[T,totalwidth=\textwidth]
  \begin{column}{0.54\textwidth}
    \begin{block}{Event ticketing platform}
      \begin{itemize}
        \item tables: \texttt{events}, \texttt{shows}, \texttt{tickets}, \texttt{orders}, \texttt{payments}
        \item нагрузка: availability, top events, order history, revenue dashboards
        \item проблема: PostgreSQL начинает дорого обслуживать read-heavy workload
      \end{itemize}
    \end{block}
  \end{column}
  \begin{column}{0.46\textwidth}
    \begin{tikzpicture}[>=Latex]
      \node[rounded corners=6pt,draw=white!20,fill=snowpanel,minimum width=4.4cm,minimum height=0.95cm] (m1) {Monolith / services};
      \node[rounded corners=6pt,draw=white!20,fill=snowpanel,minimum width=4.4cm,minimum height=0.95cm,below=0.6cm of m1] (m2) {PostgreSQL today};
      \node[rounded corners=6pt,draw=white!35,fill=snowpanelalt,minimum width=4.4cm,minimum height=0.95cm,below=0.6cm of m2] (m3) {Snowdrop target};
      \draw[white,->,thick] (m1) -- (m2);
      \draw[white,->,thick] (m2) -- (m3);
      \node[text=snowmuted,font=\small,align=center] at ($(m3.south)+(0,-0.9)$) {
        Идея не в полном rewrite,\\
        а в замене data engine\\
        для поддержанного workload subset
      };
    \end{tikzpicture}
  \end{column}
\end{columns}
\end{frame}

\begin{frame}{9. Что уже готово}
\spotlightbackground
\begin{columns}[T,totalwidth=\textwidth]
  \begin{column}{0.5\textwidth}
    \begin{block}{Технологически}
      \begin{itemize}
        \item Practical query subset
        \item Window / CTE / joins / set ops
        \item Schema fidelity subset
        \item Date/time + numeric semantics
        \item information\_schema + minimal pg\_catalog
      \end{itemize}
    \end{block}
  \end{column}
  \begin{column}{0.5\textwidth}
    \begin{block}{Migration-wise}
      \begin{itemize}
        \item \texttt{migration-preflight}
        \item \texttt{migration-run}
        \item \texttt{migration-auto}
        \item Safe translation preview
        \item Verification after load
      \end{itemize}
    \end{block}
  \end{column}
\end{columns}
\vspace{0.35cm}
\color{snowmuted}\small То есть это уже не просто research prototype, а заготовка под practical migration product for the safe subset.
\end{frame}

\begin{frame}{10. Главная мысль}
\spotlightbackground
\vspace{0.7cm}
\begin{center}
{\LARGE\bfseries Snowdrop помогает ускорять SQL workload}\\[0.25cm]
{\Large\color{snowmuted}без отказа от PostgreSQL-экосистемы и без полного архитектурного rewrite}
\end{center}

\vspace{1.0cm}
\begin{columns}[T,totalwidth=\textwidth]
  \begin{column}{0.5\textwidth}
    \begin{block}{Что уже можно говорить уверенно}
      \begin{itemize}
        \item Архитектура доказана кодом
        \item Migration path уже существует
        \item На своём целевом workload Snowdrop уже показывает ускорение
      \end{itemize}
    \end{block}
  \end{column}
  \begin{column}{0.5\textwidth}
    \begin{block}{Какой next step звучит по-взрослому}
      \begin{itemize}
        \item strengthen JOIN performance
        \item widen safe migration subset
        \item test on larger real candidate databases
      \end{itemize}
    \end{block}
  \end{column}
\end{columns}
\end{frame}

\end{document}
